\section{What it means}\label{sec:recs}

\begin{keyidea}
\textbf{Fable~5.1 at the \cref{tab:prices} prices: move the session to Sonnet, and decide once.} On Fable 5.1
at the Table~\ref{tab:prices} prices, sticky routing saved about \CfFableStickySaving\,\% (typical per-session saving: geometric mean of
per-pair cost ratios; actual spend fell \RvSpendCutAll\,\% over all \NScenarios\ scenarios, from \$\RvSpendAnchor\ to
\$\RvSpendSticky\ per session, and \RvSpendCutTest\,\% on the test split), with turn-pass
fraction inside the preregistered \NIMarginPts-point margin on the \NTest\ test scenarios. Across all
\QAllScen\ scenarios that margin was not shown for sticky (\QAllFableSticky, CI \QAllFableStickyCI), and
routed Fable arms passed the final hidden tests \FinalGapLow--\FinalGapHigh\ points less often
(exploratory). The sticky policy cost less than the shipped one (\HThreeFable\X\ on Fable, \HThreeOpus\X\ on
Opus); the preregistered test was only that sticky costs no more than shipped. At these prices the saving is
Sonnet-at-medium relative to the host at its default effort; the host at medium effort was not measured
(\cref{sec:limits}).
\end{keyidea}

\paragraph{How much of the saving is the router, and how much the effort setting.} An earlier draft read the data
as supporting ``use Sonnet'' as much as ``use the router''. That reading is withdrawn. When the sticky router chose
Sonnet it ran every main request at medium reasoning effort, while the plain-Sonnet control used the provider
default; those sticky sessions cost \RvEffFableAll\X\ (Fable) and \RvEffOpusAll\X\ (Opus) the control on the same
scenarios, with no detectable quality difference (\cref{sec:effort}). So part of the routing saving is the effort
setting, which a ``just use Sonnet'' deployment would also get if it set the same effort. On Fable the near-equality
of sticky (\CfFableStickyPair\X) and plain Sonnet (\CfFableSonnetPair\X) is two effects cancelling: medium-effort
Sonnet is cheaper than the control, while the \RvStickyHostNFable\ sticky sessions that kept the host cost
\RvStickyHostFable\X\ the plain host. The per-1,000 figures (\$\SavFableSticky\ for sticky, \$\SavFableSonnet\ for
Sonnet) are mean-difference (spend) measures. The judge's own calls are not in the session costs
(\cref{sec:design}); by an estimate from another study they would add about \JudgeShippedCents\ cents per shipped
session, which does not change these conclusions. Then a follow-up designed in advance (its plan was committed \EcFirstAfter\ before its first analysed session;
\cref{sec:effortfollow}) showed that medium effort alone makes plain Sonnet cost \EcRatio\X\ the default (95\,\% CI \EcCI),
run concurrently on the \EcScen\ test scenarios. So the effort setting alone explains most of the
sticky-versus-plain-Sonnet gap; read the other way round, sticky-on-Sonnet cost about \EcSecGapPct\,\% more than medium
plain Sonnet (\EcSecRatio; exploratory and not concurrent), within what bundle behaviour or timing could produce. The
practical reading: on Fable~5.1 at the \cref{tab:prices} prices, the saving comes from running the session on Sonnet at
medium effort. On the Sonnet sessions the router's extra machinery did not help: sticky-on-Sonnet cost about \EcSecGapPct\,\% more
than plain Sonnet at medium effort (\EcSecRatio\ the other way round; exploratory, not concurrent), and the
\RvStickyHostNFable\ sessions per host that kept the host cost \RvStickyHostFableTwo\X\ (Fable) and
\RvStickyHostOpusTwo\X\ (Opus) the plain host. This campaign shows no benefit of this router beyond the decision to use
Sonnet. Anyone already running plain Sonnet can take the effort result on its own: medium effort cut
cost by about \EcCutPct\,\% with no detectable turn-pass loss on these \EcScen\ scenarios (\RvDocsTest\ docs, \RvExplainTest\ explain, \RvReviewTest\ review); the main campaign
hints at a loss on docs, explain and review work (sticky-on-Sonnet at medium minus the default-effort control:
\RvKTDtp\ [\RvKTDtpCI] on those \RvKTN\ scenarios against \RvRestDtp\ [\RvRestDtpCI] on the rest; post hoc, and
confounded with the bundle's context) that the follow-up is too small to settle.

\begin{keyidea}
\textbf{Opus~5.5 at the \cref{tab:prices} prices: do not route to Sonnet.} Opus's cache reads are cheaper than
Sonnet's, and Sonnet needed more requests for the same scripted work, so every routing arm cost more:
\CfOpusStickyExtra--\CfOpusSonnetExtra\,\% more in the pair reading and
\CfOpusStickyArmExtra--\CfOpusSonnetArmExtra\,\% in the arm reading. The plain host is the right default
\emph{at these prices}. If Opus cache reads cost about \$\BEStickyPair\ rather than \$\PriceOpusRead\ per
million tokens, sticky routing would roughly break even (\cref{sec:breakeven}). The Opus result does not depend on
the effort setting: sticky sessions that ran Sonnet at medium effort still cost \RvEffVsAnchorOpusTwo\X\ (95\,\% CI
\RvEffVsAnchorOpusCI) the plain Opus host.
\end{keyidea}

\paragraph{Quality by task type.} The aggregate Fable sticky quality hides a split (\cref{tab:taskfable}):
\RgNTaskFail\ of \RgNTask\ task types fall below the \NIMarginPts-point margin (\RgTaskFailList), \RgNTaskThreeX\
of them by about three times, and they are among the largest savers. The aggregate passes because bugfix and
mixed scenarios are numerous. Routing only the task types that hold their quality could avoid most of the loss;
that is a hypothesis for the next campaign, not a finding.

\paragraph{The judge.} Keep Jev as the default judge. The post-hoc Cloudflare Clef judge matched its accuracy but
was \ClClefCostXLow--\ClClefCostXHigh\X\ the cost per decision and about three times slower from this machine; it is
worth re-testing if it can be hosted close to the agent or if its measured latency improves. Clef-Flash was cheaper
but less accurate (\cref{sec:clef}). Neither was part of the paired campaign. Taken together, the evidence does not
establish the value of the decision-judge machinery itself: no judge met the usefulness rule on real decisions
(\cref{sec:traces}), and the router showed no benefit beyond the decision to use Sonnet (\cref{sec:effort}). Jev
remains the best of the judges tested.

\paragraph{Caching guidance (hypotheses).} The mechanisms in \cref{sec:caching,sec:composition} suggest the
following for other hosts and prices. They are hypotheses drawn from two hosts on one provider, not tested
findings:
\begin{itemize}
  \item \textbf{Compare price per token class, not price per token.} A model that is cheaper on input and
    output may be dearer in a long conversation if its cache reads cost more.
  \item \textbf{Every switch re-writes the history.} Switching less often, ideally once at the start, should
    avoid most rebuild cost.
  \item \textbf{Count requests, not just tokens.} The cheaper model made more requests for the same scripted
    work here; that volume, not only the price, decided the Opus result.
  \item \textbf{Pauses matter.} An idle gap longer than the cache lifetime makes the next request pay to
    write the history again, on any model.
  \item \textbf{Measure with isolation.} Give every measured session a private cache key (a nonce), and never
    start two identical requests at the same instant.
\end{itemize}

\paragraph{How to use the numbers.} For a deployment that looks like this campaign, use the measured savings
per 1,000 sessions in \cref{tab:per1000,tab:per1000full}. The savings model is suitable for a pooled
forecast over a similar mix, with its interval; do not use its per-host or per-task figures
(\cref{sec:modellimits}).

\paragraph{What to measure next.}
\begin{itemize}
  \item \textbf{Longer sessions.} Real sessions often run for many dozens of turns; these ran at most
    \MaxTurns.
  \item \textbf{Real traffic.} Scripted users do not wait, wander or abandon sessions as people do.
  \item \textbf{A production measure mode.} Routed and unrouted paths cannot both run in one live session,
    but each production session can carry the model's prediction and interval, and sampled offline replays
    can provide fresh paired measurements.
  \item \textbf{Other hosts and providers.} The direction of the result depends on the price table; a new
    host model needs its own measurement, or at least a check of its read and write prices against Sonnet's.
  \item \textbf{The final-pass signal.} A larger sample would settle whether the small excess of final-test
    losses on routed Fable arms is real (\cref{sec:finalpass}).
\end{itemize}
